% Generated draft: requires mathematical and source review.
\documentclass[11pt]{article}
\usepackage[margin=25mm]{geometry}
\usepackage{fontspec}
\setmainfont{Latin Modern Roman}
\usepackage{amsmath,amssymb,mathtools,unicode-math}
\setmathfont{Latin Modern Math}
\usepackage{enumitem,needspace,xurl,hyperref}
\setlength{\parindent}{0pt}
\setlength{\parskip}{5pt}
\setlength{\emergencystretch}{3em}
\newcommand{\N}{\mathbb{N}}
\newcommand{\Z}{\mathbb{Z}}
\newcommand{\Q}{\mathbb{Q}}
\newcommand{\R}{\mathbb{R}}
\newcommand{\C}{\mathbb{C}}
\newcommand{\F}{\mathbb{F}}
\newcommand{\A}{\mathbb{A}}
\newcommand{\PP}{\mathbb{P}}
\newcommand{\E}{\mathbb{E}}
\newcommand{\eps}{\varepsilon}
\newcommand{\dd}{\,\mathrm{d}}
\begin{document}
\section*{UnsolvedMath ID: 600013}
\textbf{Source problem code:} catalogue entry
\subsection{Definitions and mathematical statement}
The source asks for conceptual proofs and possible generalizations of its projective configuration theorems, particularly the dodecagon configuration in Figure 5. Write $P=(P_1,\ldots,P_{12})$ for the outer dodecagon and $Q=(Q_1,\ldots,Q_{12})$ for the inner dodecagon produced by that figure's incidence construction. Its stated implication is
\[
\left(\exists\text{ a conic }C\quad P_i\in C\text{ for all }1\le i\le12\right)
\ \Longrightarrow\
\left(\exists\text{ a conic }D\quad Q_i\in D\text{ for all }1\le i\le12\right),
\]
for every admissible outer configuration for which the indicated construction is defined. Here $C$ and $D$ may be different conics. In the usual projective-coordinate convention, a conic is the zero locus in a projective plane of a nonzero homogeneous quadratic polynomial; whether singular conics are permitted is not supplied.

The record does not reproduce Figure 5 or give any formula or incidence rule for the $Q_i$. The implication therefore remains a schema referring to the source figure, not an independently specified theorem about arbitrary pairs of dodecagons. The underlying field, general-position conditions, allowable degeneracies, and ordering of the vertices must also be supplied from that source. No incidence construction is inferred here.

The broader request is to find conceptual proofs and generalizations of the source's configuration theorems. The supplied research summary refers to eight Schwartz--Tabachnikov configurations but does not state the other configurations or give a precise criterion for a conceptual proof or a proposed generalization.

The registry labels the item solved. Its research summary attributes a proof of the dodecagon case to Fedor Nilov, using a planar projection of a one-sheeted hyperboloid, as reported in Tabachnikov's 2016 survey, and says the proof was never published. Thus the source's original request to prove the implication and its supplied later resolution report are retained separately; no proof or publication verification is claimed here.
\subsection{Short English statement}
The requested dodecagon theorem says that the particular inner dodecagon constructed in the missing Figure 5 lies on a conic whenever the outer one does. Reconstructing that figure is essential to state the theorem independently. A supplied later summary reports Nilov's unpublished proof and labels all eight configurations solved; the broader search for conceptual proofs and generalizations remains unspecified.
\subsection{Sources}
\begin{itemize}
\item Local source registry: \url{data/unsolvedmath/problems.json}, permanent UnsolvedMath ID $600013$, source code catalogue entry, including the supplied solved classification and research summary. Dataset attribution: Ulam AI and UnsolvedMath Contributors; original-source attribution and rights are retained.
\item Tabachnikov, \emph{A Baker's Dozen of Problems} (2015), Section 11, unnumbered problems and Figure 5, as identified by the registry: \url{https://amj.math.stonybrook.edu/html-articles/Files-2015-2024/14-01/}. The metadata records HTML access on July 29, 2026, and flags current status and rights for release review.
\item Tabachnikov's 2016 survey, arXiv:1607.04758, section ``Configurations'', as cited by the supplied research summary. The summary attributes the dodecagon proof to Fedor Nilov and states that it was unpublished. No separate Nilov proof citation or complete bibliographic data for the underlying Schwartz--Tabachnikov papers is supplied.
\end{itemize}
\end{document}
